¶ 11) For any integer \(n \geq 1\) , define a map \(\varphi_n\) from \((0, \frac{1}{2} (\times \mathbf{T}^2 \text{ to } \mathbf{R}^3 \text{ by putting } \varphi_n(r; \alpha, \beta) = ((n + r \cos 2\pi\beta) \cos 2\pi\alpha, (n + r \cos 2\pi\beta) \sin 2\pi\alpha, r \sin 2\pi\beta)\) . Put

\[
\mathrm{T} _ {n} = \varphi_ {n} \left(\left[ 0, \frac {1}{2} \right[ \times \mathbf {T} ^ {2}\right), \quad \mathrm{S} _ {n} = \varphi_ {n} (\{0 \} \times \mathbf {T} ^ {2}).
\]

\begin{enumerate}
\def\labelenumi{\alph{enumi})}
\item
  Show that the restriction of \(\varphi_{n}\) to \(\left[0,\frac{1}{2}\right[\times T^{2}\) is an isomorphism (of class \(C^{\omega}\) ) from \(\left]0,\frac{1}{2}\right[\times T^{2}\) to \(T_{n}-S_{n}\) .
\item
  Let \(f_{n}:\mathrm{T}_{n}\to \mathrm{T}_{n}\) be the map that coincides with the identity on \(\mathrm{S}_n\) and is such that
\end{enumerate}

\[
f _ {n} (\varphi_ {n} (r; \alpha , \beta)) = \varphi_ {n} (r; \alpha - (n - 1) \beta , - \alpha + n \beta)
\]

for r \textgreater{} 0. Show that \(f_{n}\) induces an automorphism of \(T_{n} - S_{n}\) .

\begin{enumerate}
\def\labelenumi{\alph{enumi})}
\setcounter{enumi}{2}
\tightlist
\item
  Show that there exists a real-analytic manifold structure on \(R^{3}\) such that the maps \(f_{n}: T_{n} \to R^{3}\) and the canonical injection \(R^{3}-\bigcup_{n} S_{n} \to R^{3}\) are analytic. Denote by X the real-analytic manifold defined in this way.
\end{enumerate}

\(d)\) For \(\vartheta \in \mathbf{T}\) , denote by \(\mathrm{R}_{\vartheta}\) the rotation

\[
(x, y, z) \mapsto (x \cos 2 \pi \vartheta - y \sin 2 \pi \vartheta , x \sin 2 \pi \vartheta + y \cos 2 \pi \vartheta , z)
\]

of \(\mathbf{R}^3\) . Prove that putting, for \(\theta \in \mathbf{T}\) and \(u \in \mathrm{X}\) ,

\[
\vartheta . u = \mathrm{R} _ {\vartheta} (u) \quad \text { if } \quad u \in \mathrm{X} - \bigcup_ {n} \mathrm{S} _ {n};
\]

\(\vartheta .u = \mathrm{R}_{n\vartheta}(u)\) if \(u\in \mathrm{S}_n\) for some integer \(n\geq 1\)

defines an analytic law of operation of \(\mathbf{T}\) on \(\mathrm{X}\) .

\begin{enumerate}
\def\labelenumi{\alph{enumi})}
\setcounter{enumi}{4}
\item
  Show that the set of orbit types of X (for the operation of T defined in d)) is infinite.
\item
  Show that there is no embedding, compatible with the operation of T, of X in a finite dimensional vector space on which T operates linearly (use Exerc. 9).
\end{enumerate}

\begin{enumerate}
\def\labelenumi{\arabic{enumi})}
\setcounter{enumi}{11}
\tightlist
\item
  Let \(\mathrm{G}\) be a compact Lie group.
\end{enumerate}

\begin{enumerate}
\def\labelenumi{\alph{enumi})}
\item
  Prove that the set of conjugacy classes of normalizers of integral subgroups of G is finite (consider the operation of G on the grassmannian of subspaces of L(G), and apply Exerc. 9).
\item
  The set of conjugacy classes of (compact) semi-simple subgroups of G is finite (observe that a Lie algebra can contain only finitely-many semi-simple ideals (Chap. I, §6, Exerc. 7) and use a)).
\end{enumerate}

\begin{enumerate}
\def\labelenumi{\arabic{enumi})}
\setcounter{enumi}{12}
\item
  Let G be a Lie group, H and K two compact subgroups of G. Assume that G admits a faithful finite dimensional linear representation. Show that there exists a finite set F of subgroups of H such that for all \(g \in G\) the subgroup \(H \cap gKg^{-1}\) is conjugate in H to a subgroup in F (use Exerc. 5 and 9).
\item
  Assume that the manifold X is paracompact and locally finite dimensional. Let G be a Lie group operating properly on X, H a compact subgroup of G, and t the conjugacy class of H.
\end{enumerate}

\begin{enumerate}
\def\labelenumi{\alph{enumi})}
\item
  Show that the set \(X_{H}\) of points of X whose fixer is equal to H is a locally closed submanifold of X.
\item
  Show that the map \((g,x)\mapsto gx\) from \(G\times X_{H}\) to X induces an isomorphism (of class \(C^{r}\) ) from \(G\times^{\mathrm{N(H)}}X_{\mathrm{H}}\) to \(\mathrm{X}_{(t)}\) .
\end{enumerate}

¶ 15) Let G be a locally compact topological group operating properly on a topological space E; let \(\rho\) be a representation of G on a finite dimensional real vector space V. Denote by dg a right Haar measure on G.

\begin{enumerate}
\def\labelenumi{\alph{enumi})}
\tightlist
\item
  Let \(\mathcal{P}\) be the set of subsets A of E having the following property: there exists an open covering \((\mathrm{U}_{\alpha})_{\alpha \in \mathrm{I}}\) of E such that, for all \(\alpha \in \mathrm{I}\) , the set of \(g \in \mathrm{G}\) such that \(g\mathrm{A} \cap \mathrm{U}_{\alpha} \neq \emptyset\) is relatively compact in G.
\end{enumerate}

Prove that, for every continuous function \(f: E \to V\) whose support belongs to P, the map \(x \mapsto \int_{G} \rho(g)^{-1} f(gx) dg\) is a continuous map from E to V, compatible with the operations of G.

\begin{enumerate}
\def\labelenumi{\alph{enumi})}
\setcounter{enumi}{1}
\tightlist
\item
  Make one of the following assumptions:
\end{enumerate}

\begin{enumerate}
\def\labelenumi{(\roman{enumi})}
\item
  the space E/G is regular;
\item
  there exists an open subset U of E and a compact subset K of G such that \(\mathrm{E} = \mathrm{GU}\) and \(g\mathrm{U}\cap \mathrm{U} = \varnothing\) for \(g\notin \mathrm{K}\) .
\end{enumerate}

Show that every point x in E has a neighbourhood belonging to P (in case (i), take \(A = V \cap W\) , where V is a neighbourhood of x such that \(gV \cap V = \varnothing\) for g outside a compact subset of G, and W to be a closed neighbourhood of Gx stable under G and contained in GV).

Every point in E has an open neighbourhood stable under G and satisfying (ii).

\begin{enumerate}
\def\labelenumi{\alph{enumi})}
\setcounter{enumi}{2}
\tightlist
\item
  Assume in addition that E is completely regular. Let \(x \in E, v \in V\) be such that the fixer of x is contained in that of v. Prove that there exists a continuous map \(F : E \to V\) , compatible with the operations of G, such that \(\mathrm{F}(x) = v\) . (Let F be the space of continuous numerical functions on E whose support belongs to P, and let \(u : F \to V\) be the map \(\alpha \mapsto \int_{G} \alpha(gx) \rho(g^{-1}). v \, dg\) . Let C be a convex neighbourhood of v in V; construct a neighbourhood A of x belonging to P such that \(gx \in A\) implies that \(\rho(g^{-1}). v \in C\) , and a function \(\alpha\) on E, with support in A, such that \(\alpha(x) \neq 0\) and \(\int_{G} \alpha(gx) \, dg = 1\) . Show that \(u(\alpha)\) belongs to C and deduce that \(v \in Im u\) .)
\end{enumerate}

\begin{enumerate}
\def\labelenumi{\arabic{enumi})}
\setcounter{enumi}{15}
\tightlist
\item
  Let G be a topological group operating properly on a separated topological space E. Let x be a point in E, H its fixer in G, and S a subset of E containing x and stable under H. The group H operates on the right on \(G \times S\) by the formula \((g, s).h = (gh, h^{-1}s)\) for \(g \in G, h \in H, s \in S\) ; the map \((g, s) \mapsto gs\) induces by passage to the quotient a map from \((\mathrm{G} \times \mathrm{S})/\mathrm{H}\) to X. Then S is said to be a transversal at x if this map is a homeomorphism onto an open subset of X.
\end{enumerate}

\begin{enumerate}
\def\labelenumi{\alph{enumi})}
\item
  Show that, if \(\mathbf{G}\) is discrete, there exists a transversal at \(x\) .
\item
  Let F be a separated topological space on which G operates properly, and let \(f : E \to F\) be a continuous map, compatible with the operations of G, such that the fixer of \(f(x)\) in G is equal to H. Show that, if S is a transversal at \(f(x)\) , then \(f^{-1}(S)\) is a transversal at x.
\item
  Let N be a closed normal subgroup of G, \(\pi: E \to E/N\) the canonical projection, T a transversal at \(\pi(x)\) for the operation of G/N on E/N, and \(S \subset \pi^{-1}(T)\) a transversal at x for the operation of HN on \(\pi^{-1}(T)\) . Show that S is a transversal at x in E (for the operation of G).
\end{enumerate}

\begin{enumerate}
\def\labelenumi{\arabic{enumi})}
\setcounter{enumi}{16}
\tightlist
\item
  Let \(G\) be a Lie group. We are going to show that \(G\) has the following property:
\end{enumerate}

\begin{enumerate}
\def\labelenumi{(\Alph{enumi})}
\setcounter{enumi}{19}
\tightlist
\item
  For every completely regular topological space \(\mathrm{E}\) on which \(\mathrm{G}\) operates properly, and every point \(x\) in \(\mathrm{E}\) , there exists a transversal at \(x\) (Exerc. 16).
\end{enumerate}

\begin{enumerate}
\def\labelenumi{\alph{enumi})}
\item
  Show that every Lie group having a faithful finite dimensional linear representation satisfies (T) (use Exerc. 15 and 16 b), as well as Prop. 6).
\item
  Prove that, if G has a closed normal subgroup N such that G/N satisfies (T) and such that KN satisfies (T) for every compact subgroup K of G, then G satisfies (T) (apply Exerc. 16 c)).
\item
  If \(\mathrm{G}_0\) is compact, then G satisfies (T).
\item
  Show that, if G has a discrete normal subgroup N such that G/N satisfies (T), then G satisfies (T).
\item
  Show that G satisfies (T) (let N be the kernel of the adjoint representation; prove that \(G/N_{0}\) satisfies (T), then apply a), b) and Exerc. 9 of §1).
\end{enumerate}

\begin{enumerate}
\def\labelenumi{\arabic{enumi})}
\setcounter{enumi}{17}
\tightlist
\item
  Let \(G\) be a Lie group operating properly on a completely regular topological space \(E\) .
\end{enumerate}

\begin{enumerate}
\def\labelenumi{\alph{enumi})}
\item
  Let \(x \in E\) , and let t be its orbit type; show that there exists an open neighbourhood U of x, stable under G, such that for all \(u \in U\) the type of u is \(\geq t\) .
\item
  Assume that G operates freely on E; let \(\pi: E \to E/G\) be the canonical projection. Show that, for every point z of E/G, there exists an open neighbourhood U of z and a continuous map \(s: U \to E\) such that \(\pi \circ s(u) = u\) for all \(u \in U\) .
\end{enumerate}

\begin{enumerate}
\def\labelenumi{\arabic{enumi})}
\setcounter{enumi}{18}
\tightlist
\item
  Let G be a Lie group, H a compact subgroup of G. Prove that there exists a neighbourhood V of H such that every subgroup of G contained in V is conjugate to a subgroup of H (apply Exerc. 18 a) to the space of compact subsets of G, cf.~Integration, Chap. VIII, §5, no. 6).
\end{enumerate}

¶ 20) Let G be a compact Lie group, m a positive integer.

\begin{enumerate}
\def\labelenumi{\alph{enumi})}
\item
  Show that the set of conjugacy classes of subgroups of G of order \(\leq m\) is finite (supposing this to be false, construct a finite group F and a sequence of homomorphisms \(\varphi_{n}: F \to G\) such that \(\varphi_{i}(F)\) is not conjugate to \(\varphi_{j}(F)\) for \(i \neq j\) , and such that \(\varphi_{n}(f)\) tends to a limit \(\varphi(f)\) for all \(f \in F\) ; show that \(\varphi\) is a homomorphism, and that this leads to a contradiction with Exerc. 19).
\item
  Show that the set of conjugacy classes of subgroups F of G all of whose elements are of order \(\leq m\) is finite (let \(\mu\) be a Haar measure on G, and let U be a symmetric neighbourhood of the identity element such that \(U^{2}\) contains no non-trivial element of order \(\leq m\) ; prove that \(\operatorname{Card}(F) \leq \mu(G)/\mu(U)\) ).
\end{enumerate}

¶ 21) Let G be a compact Lie group, T a maximal torus of G.

\begin{enumerate}
\def\labelenumi{\alph{enumi})}
\item
  Let S be a set of closed subgroups of G, stable under conjugation, such that the family of subgroups \((\mathrm{S} \cap \mathrm{T})_{\mathrm{S} \in \mathcal{S}}\) is finite. Show that the set of conjugacy classes of the subgroups \(S_{0}\) , for \(S \in S\) , is finite (by using Exerc. 12 a)), reduce to the case in which the subgroups \(S_{0}\) are normal; consider the groups \(\mathrm{C}(S_{0})_{0}\) and \(\mathrm{D}(S_{0})\) , and apply Exerc. 12 b)).
\item
  Show that S is the union of a finite number of conjugacy classes of subgroups of G (by using a), reduce to the case in which the subgroups \(S \in S\) all have the same identity component \(\Sigma\) , which is normal in G; then bound the orders of the elements of the groups \(S/\Sigma\) , and apply Exerc. 20 b)).
\item
  Let E be a separated topological space on which G operates continuously. Show that if the elements of E have only a finite number of orbit types for the operation of T, the same is true for the operation of G.
\end{enumerate}

\exercisesection{Appendix~I}

\begin{enumerate}
\def\labelenumi{\arabic{enumi})}
\tightlist
\item
  Let \(G\) be a connected compact group. Denote by \(d(G)\) the upper bound of the dimensions of the quotients of \(G\) that are Lie groups. Assume that \(d(G) < \infty\) .
\end{enumerate}

\begin{enumerate}
\def\labelenumi{\alph{enumi})}
\item
  Let K be a closed normal subgroup of G; show that \(d(\mathrm{G}/\mathrm{K}) \leq d(\mathrm{G})\) , and that \(d(\mathrm{G}/\mathrm{K}) = d(\mathrm{G})\) if K is totally discontinuous.
\item
  Show that \(\mathrm{D}(\mathbf{G})\) is a Lie group, and that the kernel of the homomorphism \((x,y)\mapsto xy\) from \(\mathrm{C}(\mathrm{G})_0\times \mathrm{D}(\mathrm{G})\) to \(\mathbf{G}\) is finite.
\item
  Let \(p = d(\mathrm{C}(\mathrm{G})_{0})\) . Then \(p < \infty\) ; prove that there exists a totally discontinuous compact group D and a homomorphism \(i : Z^{p} \to D\) with dense image, such that \(\mathrm{C}(\mathrm{G})_{0}\) is isomorphic to \((\mathbf{R}^{p} \times \mathrm{D}) / \Gamma\) , where \(\Gamma\) is the image of \(Z^{p}\) under the homomorphism \(x \mapsto (x, i(x))\) (write \(\mathrm{C}(\mathrm{G})_{0}\) as a projective limit of tori of dimension p).
\item
  Assume that G is locally connected; show that G is then a Lie group.
\end{enumerate}

\specialchapter{INDEX OF NOTATION}

The reference numbers indicate respectively the chapter, paragraph and number.

\[
k \colon \mathrm{VII.Conventions}.
\]

\[
\mathrm{V} _ {\lambda} (\mathrm{S}), \mathrm{V} _ {\lambda} (s), \mathrm{V} ^ {\lambda} (\mathrm{S}), \mathrm{V} ^ {\lambda} (s) \colon \text { VII.1.1 }.
\]

\[
\mathfrak {g} _ {\lambda} (\mathfrak {h}), \mathfrak {g} ^ {\lambda} (\mathfrak {h}) \colon \text { VII.1.3 }.
\]

\[
a _ {i} (x), \operatorname{rk} (\mathfrak {g}): \text { VII.2.2 }.
\]

\[
\operatorname{Aut} (\mathfrak {g}), \operatorname{Aut} _ {e} (\mathfrak {g}): \text { VII.3.1 }.
\]

\[
\mathscr {C} ^ {\infty} \mathfrak {g} = \cap \mathscr {C} ^ {n} \mathfrak {g}: \text { VII.3.4. }
\]

\[
r _ {\rho}, r _ {\rho} ^ {0}: \text { VII.4.1 }.
\]

\[
e (\mathfrak {g}): \text {   VII.5.2.   }
\]

\[
\mathrm{A} _ {\mathrm{V}} \text {: VII.App.I.1. }
\]

\[
\mathrm{D} f \text {: VII.App.I.2.}
\]

\[
H, X _ {+}, X _ {-}: \text {   VIII.1.1.   }
\]

\[
\mathrm{V} (m): \text { VIII.1.3 }.
\]

\[
h (t), \theta (t): \text {   VIII.1.5.   }
\]

\[
\mathfrak {g}, \mathfrak {h}, \mathrm{R} (\mathfrak {g}, \mathfrak {h}) = \mathrm{R}: \text {   VIII.2.2.   }
\]

\[
\mathfrak {g} ^ {\alpha}, \mathfrak {h} _ {\alpha}, \mathfrak {s} _ {\alpha}, H _ {\alpha}, h _ {\alpha}, \theta_ {\alpha} (t), s _ {\alpha}: \text {   VIII.2.2.   }
\]

\[
\mathfrak {h} _ {\mathbf {Q}}, \mathfrak {h} _ {\mathbf {Q}} ^ {*}, \mathfrak {h} _ {\mathbf {R}}, \mathfrak {h} _ {\mathbf {R}} ^ {*}: \text {   VIII.2.2.   }
\]

\[
\mathrm{N} _ {\alpha \beta}: \mathrm{VIII.2.4}.
\]

\[
\mathfrak {g} ^ {\mathrm{P}}, \mathfrak {h} _ {\mathrm{P}}: \text {   VIII.3.1.   }
\]

\[
n (\alpha , \beta), X _ {\alpha} ^ {0}, H _ {\alpha} ^ {0}, X _ {- \alpha} ^ {0}: \text { VIII.4.2 }.
\]

\[
\theta , \mathfrak {a}, \mathfrak {a} _ {+}, \mathfrak {a} _ {-}: \text {   VIII.4.2.   }
\]

\[
x _ {\alpha \beta}, y _ {\alpha \beta}, \mathfrak {n}: \text {   VIII.4.3.   }
\]

\[
\operatorname{Aut} (\mathfrak {g}, \mathfrak {h}), \mathrm{A} (\mathrm{R}): \text {   VIII.5.1.   }
\]

\[
\mathrm{T} _ {\mathrm{P}}, \mathrm{T} _ {\mathrm{Q}}, \operatorname{Aut} _ {0} (\mathfrak {g}), \operatorname{Aut} _ {0} (\mathfrak {g}, \mathfrak {h}), \operatorname{Aut} _ {e} (\mathfrak {g}, \mathfrak {h}): \text {VIII.5.2}.
\]

\[
(\mathfrak {g}, \mathfrak {h}), \check {\mathrm{R}}, \mathrm{W}, \mathrm{B}, \mathrm{R} _ {+}, \mathrm{R} _ {-}, \mathfrak {n} _ {+}, \mathfrak {n} _ {-}, \mathfrak {b} _ {+}, \mathfrak {b} _ {-}, X _ {\alpha}, Y _ {\alpha}, H _ {\alpha}: \text {VIII.6.Conventions.}
\]

\[
w _ {0} \colon \text {   VIII.7.2.   }
\]

\[
[ \mathrm{E} ], [ \mathrm{E} ] ^ {*}, \mathscr {R} (\mathfrak {a}): \text {   VIII.7.6.   }
\]

\[
\mathbf {Z} [ \Delta ], e ^ {\lambda}, \operatorname{ch} (\mathrm{E}) \colon \text { VIII.7.7 }.
\]

\[
\exp (\lambda): \text {   VIII.8.1.   }
\]

\[
\theta (a), \theta^ {*} (a), \mathrm{I} (\mathfrak {g}), \mathrm{I} (\mathfrak {g} ^ {*}), \mathrm{S} (\mathfrak {g}), \mathrm{S} (\mathfrak {g} ^ {*}) \colon \text { VIII.8.3 }.
\]

\[
u ^ {\sharp}, \chi_ {\lambda}, \eta , \delta , \theta , \varphi : \mathrm{VIII.8.5}.
\]

\[
\mathfrak {B} (\nu), \mathrm{K}, d: \text {   VIII.9.1.   }
\]

\[
\mathrm{J} (e ^ {\mu}) \colon \text { VIII.9.2. }
\]

\((x,h,y)\) : VIII.11.1. \(x^{(n,d)},x^{(n)},\binom{x}{n}\) : VIII.12.2. \(\mathbf{x}^{(n)}\) , \([n],c_r,\mathcal{G}\) : VIII.12.6. \(\mathcal{H},\mathcal{U}_{+},\mathcal{U}_{-},\mathcal{U}_{0},\mathcal{L},\mathcal{G}\) : VIII.12.6. \(\mathfrak{o}(\Psi)\) : VIII.13.2. \(\mathfrak{sp}(\Psi)\) : VIII.13.3. \(\mathrm{G}_0,\mathrm{C}(\mathrm{G}),\mathrm{D}(\mathrm{G})\) : IX.1.Conventions. \(\mathrm{N_G(H),Z_G(H),N(H),Z(H)}\) : IX.1.Conventions. \(\Gamma (\mathrm{G})\) : IX.1.2. \(\tilde{\mathrm{D}} (\mathrm{G})\) : IX.2.3.\\
rk G: IX.2.4. \(\mathrm{W_G(T),W(T)}\) : IX.2.5. \(\mathfrak{a}_{[\mathbb{R}]},\mathfrak{g}_{(\mathbb{C})},\mathfrak{g}_{\mathbb{C}}\) : IX.3.1. \(\mathfrak{a}_u,u_\alpha ,v_\alpha\) : IX.3.2. \(\mathbf{U}(\varPhi),\mathbf{SU}(\varPhi),\mathfrak{u}(\varPhi),\mathfrak{s}\mathfrak{u}(\varPhi)\) : IX.3.4. \(\mathbf{U}(n,\mathbf{C}),\mathbf{SU}(n,\mathbf{C}),\mathfrak{u}(n,\mathbf{C}),\mathfrak{s}\mathfrak{u}(n,\mathbf{C})\) : IX.3.4. \(\mathbf{O}(\mathrm{Q}),\mathbf{SO}(\mathrm{Q}),\mathfrak{o}(\mathrm{Q})\) : IX.3.5. \(\mathbf{O}(n,\mathbf{R}),\mathbf{SO}(n,\mathbf{R}),\mathfrak{o}(n,\mathbf{R})\) : IX.3.5. \(U,V,K,\theta\) : IX.3.6.\\
X(H),X(f): IX.4.1. \(\delta (a)\) : IX.4.1. \(\Gamma (f)\) : IX.4.2. \(\langle a,X\rangle\) : IX.4.2. \(\tilde{\rho},\tilde{\mathrm{V}},\tilde{\mathrm{V}}_{\lambda}(\mathrm{G})\) : IX.4.3.\\
P(ρ,T): IX.4.3. \(\tilde{\mathrm{L}} (\rho)\) : IX.4.3.\\
R(G,T): IX.4.4.\\
\(\mathrm{V}(\alpha),\mathrm{K}(\alpha),\iota_{\alpha},\rho_{\alpha}\): IX.4.4.\\
\(\mathrm{Z}_{\alpha},\mathrm{S}_{\alpha}\): IX.4.5.\\
\(\mathrm{K}_{\alpha}\): IX.4.5.\\
\(\mathrm{R}^{\vee}(\mathrm{G},\mathrm{T})\): IX.4.5.\\
N(G,T): IX.4.6.\\
D*(G,T),D*(G,T): IX.4.9.\\
D*(f),D*(f): IX.4.9.\\
D*(G),D*(G): IX.4.9.\\
Aut(G),Aut(G,T): IX.4.10.\\
\(\mathrm{G}_{r},\mathrm{T}_{r}\): IX.5.1.\\
\(\mathrm{W}_{\alpha},\mathrm{W}'_{\alpha}\): IX.5.2.\\
\(\mathrm{H}_{\alpha,n},t_{r}\): IX.5.2.\\
\(\mathrm{H}_{A}\): IX.5.2.\\
\(\mathfrak{g}_{\mathrm{reg}}\): IX.5.2.\\
\(f,f_{r}\): IX.5.4.\\
\(\varphi,\varphi_{r},\varphi_{A}\): IX.5.4.\\
\(\mathfrak{g},r,\psi_{r}\): IX.5.4

\(w(\mathrm{G})\) : IX.6.Conventions. \(u\cap v\) : IX.6.1. \(\omega_{\mathrm{G / T}}\cap \omega_{\mathrm{T}}\) : IX.6.2. \(\mathrm{Ad}_{\mathfrak{g} / \mathfrak{t}}(t),\delta_{\mathrm{G}}(t)\) : IX.6.2. \(\mathrm{R}_+(\mathrm{G},\mathrm{T})\) : IX.6.2. \(\lambda_{\mathfrak{h}}\) : IX.6.3. \(\pi_{\mathfrak{g}}\) : IX.6.3. \(\mathcal{S}^r (\mathrm{X};\mathrm{E})\) : IX.6.4. \(s^\sharp\) : IX.6.4. \(\Omega (\mathrm{X}),\Omega (\mathrm{X})^{\mathrm{G}}\) : IX.6.5. \(\Gamma (\mathrm{T})_{++},\mathrm{X}_{+},\mathrm{R}_{+},\mathrm{R}_{-}\) : IX.7.Conventions. \(\mathrm{X}_{++}\) : IX.7.1. \(\varpi_{\alpha},\rho\) : IX.7.1. \(\Theta (\mathrm{G})\) : IX.7.2. \(\operatorname {R}(\operatorname {G})\) : IX.7.3. \([\tau ],\operatorname {Ch}(\tau)\) : IX.7.3. \(\mathrm{Z}\Theta (\mathrm{G})\) : IX.7.3. \(^w f,\mathrm{J}(f)\) : IX.7.4. \(\chi_{\lambda}\) : IX.7.4. \(\mathrm{ZL}^2 (\mathrm{G})\) : IX.7.4. \(\mathrm{Q}_{\Gamma},\tilde{\Gamma} (\tau)\) : IX.7.6. \(\hat{\mathbf{G}}\) : IX.8.1. \(\operatorname {E}_u,d(u),\| \mathrm{A}\|_\infty ,\| \mathrm{A}\| _2,\langle \mathrm{A}|\mathrm{B}\rangle\) : IX.8.1.\\
F( \(\hat{\mathrm{G}}\) ), \(\mathrm{L}^2 (\mathrm{G})\) : IX.8.1. \(u(f),\overline{\mathcal{F}} (f)\) : IX.8.1. \(\mathcal{F}_u\mathrm{A},\mathcal{F}(\mathrm{A})\) : IX.8.1. \(\gamma (s)f,\delta (s)f\) : IX.8.1. \(\mathrm{L}_t f,\mathrm{R}_t f\) : IX.8.2. \(\lambda (u),\tilde{\Gamma} (u)\) : IX.8.2. \(\varphi \preccurlyeq \psi\) : IX.8.2. \(\mathcal{S}(\hat{\mathrm{G}})\) : IX.8.2. \(\chi_u\) : IX.8.3. \(\mathrm{Z}\mathcal{C}^{\infty}(\mathrm{G}),\mathrm{Z}\mathcal{C}(\mathrm{G})\) : IX.8.4. \(\operatorname {E}_{(t)}\) : IX.9.4.\\
Spin(n,R):IX.3,Exerc.5.

\specialchapter{INDEX OF TERMINOLOGY}

The reference numbers indicate respectively the chapter, paragraph and number.

absolutely simple (Lie algebra): VIII.3.2 adjoint (representation of SL(2, k)): VIII.1.4 admissible (lattice): VIII.12.8 alcove: IX.5.2 anti-hermitian (endomorphism): IX.1.1 anti-invariant: IX.7.4 biorder (of a Lie algebra): VIII.12.1 Borel (subalgebra): VIII.3.3, VIII.3.5 canonical (Cartan subalgebra, root system, basis): VIII.5.3 canonical (involution of sl(2, k)): VIII.1.1 Cartan (subalgebra): VII.2.1 Casimir element: IX.7.6 central (character of a g-module): VIII.6.1 central (function): IX.8.3 chamber: IX.5.2 character (of a g-module): VIII.7.7 Chevalley (order): VIII.12.7 Chevalley (system): VIII.2.4 classical (splittable Lie algebra): VIII.3.2 clean (subgroup): IX.4, Exerc. 15 Clebsch-Gordan (formula): VIII.9.4 compact (Lie algebra): IX.1.3 compact (real form of a complex Lie algebra): IX.3.1 compatible (representations): VII.3.1, VIII.1.4 condition (AC): VII.1.1 conjugation: IX.3.1 contravariant diagram: IX.4.9 covariant diagram: IX.4.9 Coxeter element: IX.4, Exerc. 14 dominant: IX.7.1 eigenvalue: VII.1.1 eigenvector: VII.1.1

elementary (automorphism): VII.3.1 embedding (of a manifold in the neighbourhood of a subset): IX.9.1 exceptional (splittable simple Lie algebra): VIII.3.2 exponential: VIII.8.1 facet (associated to a parabolic subset): VIII.3.4 Fitting (decomposition of a module): VII.1.1 flag: VIII.13.1 Fourier cotransform: IX.8.1 Fourier transform: IX.8.1 framing: IX.4.10 framing (of a split semi-simple Lie algebra): VIII.4.1 framed (semi-simple Lie algebra): VIII.4.1 fundamental (dominant weight): IX.7.1 fundamental (representations, simple g-modules): VIII.7.2 generating family (defined by a framing): VIII.4.1 Harish-Chandra (homomorphism): VIII.6.4 Hermann Weyl (formula of --): VIII.9.1 homogeneous symmetric space: IX.1, Exerc. 8 invariant (polynomial function): VIII.8.3 involutive (Lie algebra): IX.1, Exerc. 7 irreducible (involutive Lie algebra): IX.1, Exerc. 7 irreducible (homogeneous symmetric space): IX.1, Exerc. 8 isotropic (flag): VIII.13.2, VIII.13.3 isotypical (component of highest weight λ): VIII.7.2 Jacobson-Morozov (theorem): VIII.11.2 lattice (in a Q-vector space): VIII.12.1 linear tube: IX.9.3 maximal torus: IX.2.2 minuscule (weight): VIII.7.3 moderately increasing: IX.8.2 nilpotent (component of an element): VII.1.3 nilspace: VII.1.1 nodal group (of a torus): IX.4.2 nodal vector: IX.4.5 order (of a Q-algebra): VIII.12.1 orthogonal (Lie algebra): VIII.13.2 orthogonal (irreducible representation): VIII.7.5 parabolic (subalgebra): VIII.3.4, VIII.3.5 partition (into positive roots): VIII.9.1 permissible (lattice): VIII.12.6 positive root: IX.7.1 primary (subspace): VII.1.1 primitive (element of a module): VIII.1.2, VIII,6.1 principal (nilpotent element, simple element, sl₂-triplet): VIII.11.4 principal (orbit): IX.9.4 principal (subgroup): IX.4, Exerc. 18 proper (subspace): VII.1.1 quasi-maximal (isotropic flag): VIII.13.4 R-extremal (element): VIII.7.2 R-saturated (subset): VIII.7.2 radical (subgroup): IX.4.7 rank (of a Lie algebra): VII.2.2 rank (of a Lie group): IX.2.4 rapidly decreasing: IX.8.2 real form (of a complex Lie algebra): IX.3.1 reduced (root diagram): IX.4.8 regular (element for a linear representation): VII.4.1 regular (element of a Lie algebra): VII.2.2 regular (element of a Lie group): VII.4.2 representation ring (of a Lie algebra): VIII.7.6 root: IX.4.4 root (of (g,h)): VIII.2.2 root diagram: IX.4.8 semi-simple (component of an element): VII.1.3 semi-spinorial (representation): VIII.13.4 simple (element): VIII.11.3 simple (root): IX.7.1 singular hyperplane: IX.5.2 sl₂-triplet: VIII.11.1 spinorial (representation): VIII.13.2, VIII.13.4 split (reductive Lie algebra): VIII.2.1 splittable (Lie subalgebra): VII.5.1, VIII.10 splittable (Cartan subalgebra, reductive Lie algebra): VIII.2.1 splittable envelope (of a subalgebra of gl(V)): VII.5.2 subgroup (Cartan --): IX.2.2 subgroup (of maximum rank): IX.2.4 subgroup (principal --): IX.4, Exerc. 18 symplectic (Lie algebra): VIII.13.3 symplectic (irreducible representation): VIII.7.5 tangent (linear map of a polynomial map): VII.App.I.2 topology of compact Cʳ-convergence: IX.6.4 torsion prime number: IX.5, Exerc. 9 to 11 torus: IX.1.2 transversal: IX.9.3 and IX.9, Exerc. 16 type (An,Bn,\ldots): IX.3.3 type (of a representation): IX.App.II.1, IX.App.II.2 vector (group): IX.1.2 very regular (element): IX.5.1

weights (of a g-module): VII.1.1, VIII.1.2, VIII.6.1

weights (of a representation): IX.4.3

weights (group of - of a split Lie algebra): VIII.2.2

Weyl (group): IX.2.5

Weyl (group of a split algebra): VIII.2.2

Witt (basis): VIII.13.2, VIII.13.3, VIII.13.4

Zariski (topology): VII.App.I.1